%% =============================================================================
%% PAPER 1 - Appendix
%% =============================================================================


%% ---- A.1 BTP Italia: Instrument Details and Arbitrage Mechanics ----
\subsection{BTP Italia: Instrument Details and Arbitrage Mechanics}
\label{app:btp_details}

\subsubsection{Inflation Floor and Semiannual Principal Revaluation}
\label{app:btp_floor}
% BTP Italia: inflation-linked bond issued by Italian Treasury
% Semiannual revaluation of principal (not at maturity like BTPei)
% Embedded inflation floor at 0% on each coupon period
% Eliminates ballooning risk (unlike TIPS or BTPei)

\subsubsection{Repo and Collateral Treatment}
\label{app:btp_repo}
% BTP Italia accepted as general collateral (same as nominal BTPs)
% Repo rate identical -> cannot explain basis via funding cost differential
% Contrast with TIPS (different repo market, higher haircuts)

\subsubsection{Comparison with US TIPS}
\label{app:btp_vs_tips}
% Fleckenstein et al. (2014): TIPS basis averaged 60bps, peaked 200bps
% Today: TIPS basis ~0bps (QE, broader participation)
% Italian IL basis persists >10 years
% Different investor base: retail (BTP Italia) vs institutional (TIPS)


%% ---- A.2 Trade Construction Details ----
\subsection{Trade Construction Details}
\label{app:trade_details}

\subsubsection{Entry/Exit Rules and Thresholds}
\label{app:entry_exit}
Table~\ref{tab:entry_exit_rules} summarises entry and exit thresholds
for each strategy. All strategies require a minimum of 6 months to
maturity at entry and maintain at least one open trade after the first
signal. The CDS--Bond universe is restricted to single-name CDS whose
reference entity belongs to the current on-the-run iTraxx Main series;
names that exit the index are no longer traded. All CDS are centrally
cleared. The iTraxx skew strategy trades only on-the-run series across
Main, SnrFin, SubFin, and Xover.

\input{\tablepath/entry_exit_thresholds_article}

\subsubsection{P\&L Calculation}
\label{app:pnl}
% P&L = Capital Gain + Carry
% Capital Gain = sign * (prev_basis - current_basis) * DV01
% Carry = sign * prev_basis * (days/365)
% DV01 updated daily (current, not entry)

\subsubsection{Trade-Level Statistics}
\label{app:trade_stats}

\input{\tablepath/trade_statistics_article}

\noindent\textit{A.2.4\quad Transaction Costs}
\label{app:fees}
Transaction costs are computed at the individual trade level and applied
as entry and exit fees proportional to DV01 (bonds) or duration (CDS).
Fee levels are regime-dependent: at each trade date, the prevailing level
of the iTraxx Main 5Y index determines which of three tiers applies
(Table~\ref{tab:fee_schedule}). Thresholds are set at 60 and 100~bps,
reflecting the empirical widening of bid-ask spreads in stressed markets.
For trades held to maturity, the exit fee is waived. All return series
in Sections~3--6 are net of these costs.

\input{\tablepath/fee_schedule_article}

%% ---- A.3 Signal-Weighted vs Equal-Weighted Comparison ----
\subsection{Signal-Weighted vs Equal-Weighted Comparison}
\label{app:sw_ew}

Table~\ref{tab:sw_ew_basic} compares the signal-weighted (SW,
\citealp{rebonato2021convexity}) and equal-weighted (EW,
\citealp{duarte2007risk}) indices. SW is adopted as the baseline in the
main text; EW results confirm robustness to the weighting scheme.

\input{\tablepath/SW_EW_basic_metrics_3strategies_article}


%% ---- A.4 Candidate Factor Universe ----
\subsection{Candidate Factor Universe}
\label{app:factors}

\subsubsection{Complete Factor List with Sources}
\label{app:factor_list}
Table~\ref{tab:factor_list} reports the full set of candidate risk
factors used in the penalised regressions of Section~5, together with
their data sources, transformations, and economic category.
\input{\tablepath/factor_list}

\subsubsection{Stationarity Tests}
\label{app:stationarity}
% ADF + KPSS on all factors
% 63 stationary, 3 differenced (CREDIT_EU, YSP_US, YSP_EU)
% 6 borderline discussed
% TABLE REF: \ref{tab:stationarity}


%% ---- A.5 PCA Robustness ----
\subsection{PCA Robustness}
\label{app:pca_robustness}

\subsubsection{Robustness across Window Length and Number of Components}
\label{app:pca_wk}

Table~\ref{tab:pca_robustness} reports alpha and $\bar{R}^2$ for all
$(W,K)$ combinations. Results are robust to the choice of window length
and number of components.

\input{../results/robustness/tables/PCA_robustness_alpha_R2_article}

\subsubsection{Subperiod and Regime Analysis}
\label{app:pca_subperiod}

Half-sample splits and regime-dependent alpha confirm that PCA-based
alpha is stable across sub-periods and increases during stress.

\input{\tablepath/PCA_subperiod_results_article}

\input{\tablepath/PCA_regime_summary_article}

\begin{figure}[H]
\centering
\includegraphics[width=\textwidth]{\figpath/pca/pca_rolling_alpha.pdf}
\caption{Rolling Alpha with Regime Shading (PCA Factors).
Grey bands: 95\% confidence intervals.
Red shading: iTraxx Main 5Y $> 100$ bps.}
\label{fig:app_pca_rolling_alpha}
\end{figure}

\subsubsection{Conditional Alpha: Ferson--Schadt and Threshold Robustness (PCA)}
\label{app:pca_conditional}

Table~\ref{tab:pca_cond_fs_contemporaneous_article} reports Ferson--Schadt
conditional alpha using PCA factors.
Table~\ref{tab:pca_cond_threshold_contemporaneous_article} confirms
robustness across stress thresholds (80, 100, 120 bps).

\input{\tablepath/PCA_Cond_Alpha_FS_article_contemporaneous}

\input{\tablepath/PCA_Cond_Alpha_Threshold_article_contemporaneous}

\subsubsection{Variance and Loadings Diagnostics}
\label{app:pca_diagnostics}

Figure~\ref{fig:app_pca_variance_ts} shows the time-varying cumulative
variance explained.

\begin{figure}[H]
\centering
\includegraphics[width=0.85\textwidth]{\figpath/pca/pca_variance_timeseries.pdf}
\caption{Time Series of Cumulative Variance Explained by Rolling PCA.}
\label{fig:app_pca_variance_ts}
\end{figure}




%% ---- A.6 AEN Robustness ----
\subsection{AEN Robustness}
\label{app:aen_robustness}

\subsubsection{Sensitivity to Hyperparameters}
\label{app:aen_sensitivity}

We verify that factor selection, model dimension, and post-selection alpha are invariant to the correlation pre-cleaning threshold ($\rho \in \{0.75, 0.80, 0.85, 0.90, 0.95\}$) and the adaptive weight exponent ($\gamma \in \{0.5, 1, 2, 3\}$).
In all configurations, the AEN selects the same factors with identical post-selection OLS results.
No selected factor belongs to a highly correlated pair, so varying the cleaning criterion does not alter the surviving set; similarly, the Stage~1 elastic net coefficients exhibit sufficient separation between relevant and irrelevant factors to make the adaptive weighting insensitive to the exponent.
Detailed results are available upon request.

\subsubsection{Factor Selection Across Methods}
\label{app:aen_selection_matrix}

Table~\ref{tab:aen_selection_names} reports the factors selected by each penalised method for each strategy.
LASSO, Elastic Net, and Adaptive LASSO produce broadly overlapping selections, while AEN---which retains only factors exceeding the 60\% bootstrap stability threshold---yields the sparsest models.
ALASSO-Chen selects substantially more factors across all strategies, consistent with the lighter penalisation implied by its data-driven $\lambda_1$ rule \citep{chen2008extended}.

\input{\tablepath/AEN_Selection_Names_article}


\subsubsection{Conditional Alpha: Detailed Results (AEN)}
\label{app:aen_conditional}

Table~\ref{tab:conditional_alpha_fs} reports the full Ferson--Schadt
conditional alpha model with AEN-selected factors.
Table~\ref{tab:conditional_alpha_threshold} confirms robustness across
stress thresholds (80, 100, 120 bps).

\input{../results/aen/hqc/tables/Cond_Alpha_FS_Thesis}

\input{../results/aen/hqc/tables/Cond_Alpha_Threshold_Thesis}


\subsubsection{Rolling Alpha}
\label{app:aen_rolling}

Figure~\ref{fig:app_aen_rolling_alpha}
plot rolling 36-month alpha with regime shading for AEN and PCA factors.
Alpha spikes during stress episodes under both methods.

\begin{figure}[H]
\centering
\includegraphics[width=\textwidth]{\figpath/aen/aen_rolling_alpha.pdf}
\caption{Rolling Alpha with Regime Shading (AEN Stable Factors).
Grey bands: 95\% confidence intervals.
Red shading: iTraxx Main 5Y $> 100$ bps. Rolling window: 36 months.}
\label{fig:app_aen_rolling_alpha}
\end{figure}


%% ---- A.7 Subperiod and Rolling Window Analysis ----
\subsection{Subperiod and Rolling Window Analysis}
\label{app:subperiod}
 
This appendix reports sub-period splits and threshold robustness
for the benchmark factor models of Section~4. Each table contains
Panel~A (full sample, first half, second half, HIGH, and NORMAL regimes)
and Panel~B (stress threshold robustness at 80, 100, and 120 bps of
iTraxx Main 5Y).
 
 
%% ---- A.7.1 Duarte et al. (2007) ----
\subsubsection{Duarte et al.\ (2007)}
\label{app:subperiod_duarte}
Table~\ref{tab:subperiod_regime_duarte} reports sub-period splits and stress
regime alpha for the \citet{duarte2007risk} specification.
\input{\tablepath/subperiod_regime_duarte_monthly}
 
\begin{figure}[H]
\centering
\includegraphics[width=\textwidth]{\figpath/rolling_alpha_regime_btp_italia_duarte_monthly.pdf}
\includegraphics[width=\textwidth]{\figpath/rolling_alpha_regime_cds_bond_basis_duarte_monthly.pdf}
\includegraphics[width=\textwidth]{\figpath/rolling_alpha_regime_itraxx_combined_duarte_monthly.pdf}
\caption{Rolling 36-Month Alpha with Regime Shading --- Duarte et al.\ (2007).
  95\% confidence bands. Red: HIGH stress (iTraxx Main $> 100$ bps).}
\label{fig:app_rolling_duarte}
\end{figure}

%% ---- A.7.2 Brooks et al. (2020) / Active FI ----
\subsubsection{Brooks et al.\ (2020) --- Active Fixed Income}
\label{app:subperiod_activefi}
Table~\ref{tab:subperiod_regime_activefi} reports the same analysis under the
\citet{brooks2020active} specification.
\input{\tablepath/subperiod_regime_activefi_monthly}
 
\begin{figure}[H]
\centering
\includegraphics[width=\textwidth]{\figpath/rolling_alpha_regime_btp_italia_activefi_monthly.pdf}
\includegraphics[width=\textwidth]{\figpath/rolling_alpha_regime_cds_bond_basis_activefi_monthly.pdf}
\includegraphics[width=\textwidth]{\figpath/rolling_alpha_regime_itraxx_combined_activefi_monthly.pdf}
\caption{Rolling 36-Month Alpha with Regime Shading --- Brooks et al.\ (2020).
  Same format as Figure~\ref{fig:app_rolling_duarte}.}
\label{fig:app_rolling_activefi}
\end{figure}
 
 
%% ---- A.7.3 Fung & Hsieh (2004) ----
\subsubsection{Fung \& Hsieh (2004)}
\label{app:subperiod_fh}
Table~\ref{tab:subperiod_regime_fh} reports results for the
\citet{fung2002risk} specification.

 
\input{\tablepath/subperiod_regime_fh_monthly}
 
\begin{figure}[H]
\centering
\includegraphics[width=\textwidth]{\figpath/rolling_alpha_regime_btp_italia_funghsieh_monthly.pdf}
\includegraphics[width=\textwidth]{\figpath/rolling_alpha_regime_cds_bond_basis_funghsieh_monthly.pdf}
\includegraphics[width=\textwidth]{\figpath/rolling_alpha_regime_itraxx_combined_funghsieh_monthly.pdf}
\caption{Rolling 36-Month Alpha with Regime Shading --- Fung \& Hsieh (2004).
  Same format as Figure~\ref{fig:app_rolling_duarte}.}
\label{fig:app_rolling_fh}
\end{figure}


%% ---- A.8 Cross-Strategy Interdependencies: Additional Tests ----
\subsection{Cross-Strategy Interdependencies: Additional Tests}
\label{app:rq3_additional}
 
This appendix collects robustness checks and supplementary analyses for
Section~\ref{sec:interdependencies}. All tests use the same mispricing
magnitudes $m_{i,t}$ and regime definitions as the main text.
 
 
%% ---- A.8.1 Purged Correlations ----
\subsubsection{Purged Correlations}
\label{app:purged_corr}
 
To verify that cross-strategy co-movement is not driven solely by common
risk-factor exposures, we regress each $\Delta m_i$ on the AEN-selected
factors from Section~5 and compute correlations on the residuals.
Table~\ref{tab:rq3_purged_corr} shows that pairwise correlations remain
positive and significant after purging, indicating that the interdependencies
documented in Section~6.1 reflect capital-market linkages beyond shared
factor exposures.
 
\input{\tablepath/RQ3_Purged_Correlations_article}
 
 
%% ---- A.8.2 DCC-GARCH ----
\subsubsection{DCC-GARCH}
\label{app:dcc}
 
We estimate bivariate DCC-GARCH(1,1) models on each pair
of $\Delta m$ series. Figure~\ref{fig:app_dcc} plots the time-varying
conditional correlations. Spikes in conditional correlation coincide with
stress episodes (COVID-19, ECB hiking cycle), corroborating the
regime-dependent results in Section~6.1.
 
\begin{figure}[H]
\centering
\includegraphics[width=\textwidth]{../results/rq3_duffie/figures/F1_dcc_garch.pdf}
\caption{Dynamic Conditional Correlations (DCC-GARCH).
  Bivariate DCC(1,1) estimates for each pair of $\Delta m$.
  Shaded areas: HIGH stress regime (iTraxx Main $> 100$ bps).}
\label{fig:app_dcc}
\end{figure}
 
 

 
%% ---- A.8.6 Granger Causality: Lag Robustness ----
\subsubsection{Granger Causality: Lag Robustness}
\label{app:granger_multilag}

Table~\ref{tab:rq3_granger_multilag} reports Granger causality tests
at lag lengths $p = 1, 2, 3$. The main result --- iTraxx Combined
Granger-causes BTP Italia --- is robust across all specifications.

\input{\tablepath/RQ3_Granger_Multilag_article}

 
%% ---- A.8.7 Alternative Stress Proxies ----
\subsubsection{Alternative Stress Proxies}
\label{app:alt_stress}
 
Table~\ref{tab:rq3_alt_proxies} replicates key tests from Section~6
(unconditional correlations, regime correlations, interaction regressions)
using iTraxx Crossover, V2X, and VIX as alternative conditioning variables
instead of iTraxx Main. Results are qualitatively unchanged, confirming that
the documented interdependencies are not an artifact of the stress-proxy choice.
 
\input{\tablepath/RQ3_Alternative_Proxies_article}
 
 
%% ---- A.8.8 Threshold Clustering ----
\subsubsection{Threshold Clustering}
\label{app:clustering}
 

 
 
%% ---- A.8.9 Pairwise Long-Sample Analysis ----
\subsubsection{Pairwise Long-Sample Analysis}
\label{app:pairwise}
 

%% ---- A.9 VIF Diagnostics ----
\subsection{VIF Diagnostics}
\label{app:vif}
 
Tables~\ref{tab:vif_duarte}--\ref{tab:vif_funghsieh} report Variance
Inflation Factors for the EUR factor specifications used in the benchmark
regressions of Section~4. VIF is computed on the factor matrix for each
strategy's sample period; because the factors are identical across strategies,
differences in VIF reflect only the different estimation windows.
 
All equity and bond-market factors exhibit VIF below~10, confirming
acceptable multicollinearity levels. The Duarte specification shows
elevated VIF for the term-structure factors R2, R5, and R10
(VIF~$\approx$~15--55), which is expected given the high correlation
structure of yield-curve returns and does not affect inference on alpha.
The Active~FI and Fung~\&~Hsieh specifications show VIF below~5 for
all factors.
 
\input{\tablepath/Duarte_VIF_article_monthly}
 
\input{\tablepath/ActiveFI_VIF_article_monthly}
 
\input{\tablepath/FungHsieh_VIF_article_monthly}